

%--------------------------------------
%	PACKAGES AND OTHER DOCUMENT CONFIGURATIONS
%--------------------------------------


\documentclass[12pt, a4paper, sans]{moderncv} 
% Font sizes: 10, 11, or 12
% Paper sizes: a4paper, letterpaper, a5paper, legalpaper, executivepaper or landscape
% Font families: sans or roman

\moderncvstyle{classic} 
% CV theme - options include: 'casual' (default), 'classic', 'oldstyle' and 'banking'

\moderncvcolor{blue} 
% CV color - options include: 'blue' (default), 'orange', 'green', 'red', 'purple', 'grey' and 'black'

\usepackage{multicol}

\usepackage{lipsum} 
% Used for inserting dummy 'Lorem ipsum' text into the template

\usepackage[hyperref]{}

\usepackage[scale=0.9]{geometry} 
% Reduce document margins

%\setlength{\hintscolumnwidth}{3cm} 
% Uncomment to change the width of the dates column

%\setlength{\makecvtitlenamewidth}{10cm} 
% For the 'classic' style, uncomment to adjust the width of the space allocated to your name

\usepackage[T2A]{fontenc}
\usepackage[utf8]{inputenc}
\usepackage[russian]{babel}
%Russian-specific packages


%--------------------------------------
%	NAME AND CONTACT INFORMATION SECTION
%--------------------------------------


\firstname{Захаров Игорь \newline} 
\familyname{Александрович} 

\address{}{Москва, Россия}
\mobile{+7 (916) 624 94 84}
\email{ia2002zakharov@gmail.com}
\extrainfo{\href{https://t.me/IgorZzzzzz}{Telegram: @IgorZzzzzz}}

\photo[100pt][0.4pt]{pictures/photo} 
% The first bracket is the picture height, the second is the thickness of the frame around the picture (0pt for no frame)


%--------------------------------------
%	CURRICULUM VITAE
%--------------------------------------


\begin{document}

\makecvtitle 
% Print the CV title


%--------------------------------------
%	COMPUTER SKILLS SECTION
%--------------------------------------


\section{Skills}

\cvitem{Языки пр.}{Python, C++, C}
\cvitem{Навыки}{SQL, Git, Bash, Docker, Airflow, TeX}
\cvitem{Языки}{Русский, Английский С1}


%--------------------------------------
%	EXPERIENCE SECTION
%--------------------------------------


\section{Experience}

\cventry{Янв. 2023 -- н. в.}{ООО ВБ Тех}{\newline Аналитик в команде ранжирования.}{}{}{\begin{itemize}
\item Учет цены в ранжировании.
\item Разработка сервиса для мониторинга изменений коэффициентов ранжирования.
\item Telegram-бот для уведомлений об ошибках в airflow.
\end{itemize}}


%--------------------------------------
%	EDUCATION SECTION
%--------------------------------------


\section{Education}

\cventry{Сент. 2024 -- н. в.}{Магистратура НИУ ВШЭ}{\newline Факультет компьютерных наук,\newline образовательная программа «Магистр по наукам о данных»}{}{}{}

\cventry{Сент. 2020 -- Июнь 2024}{Бакалавриант НИУ ВШЭ}{\newline Факультет компьютерных наук,\newline образовательная программа «Прикладная математика и информатика»,\newline специализация «Анализ данных и интеллектуальные системы»}{}{}{}

\cventry{Сент. 2016 -- Май 2020}{ГБОУ Школа №57}{\newline Математический класс}{}{}{}

%--------------------------------------
%	PROJECTS SECTION
%--------------------------------------


\section{Projects}

\cventry{Окт. 2021 -- Май 2022}{Разработка прототипа пошаговой стратегической игры}{\textit{\newlineЧумаков Олег Олегович, CEO Luden.io}}{}{}{
\begin{itemize}
\item Процедурная генерация игрового поля и игрового окружения.
\item Игровой логики и взаимодействие персонажей.
\item Контроллер передвижения камеры.
\item Участие в "Game Off 2021 jam".
\item \href{https://arthadanel.itch.io/frantic-bugs}{Страница игры на itch.io.}
\end{itemize}}


%--------------------------------------
%	AWARDS SECTION
%--------------------------------------


\newpage
\section{Achievements}

\cventry{2020}{Всероссийская олимпиада школьников по программированию}{}{}{}{
\begin{itemize}
\item Призер
\end{itemize}}

\cventry{2020}{Всесибирская олимпиада школьников по программированию}{}{}{}{
\begin{itemize}
\item Диплом I степени
\end{itemize}}

\cventry{2019}{Всероссийская командная олимпиада школьников по программированию (ВКОШП)}{}{}{}{
\begin{itemize}
\item Призер
\end{itemize}}

\cventry{2018}{Московская олимпиада школьников по информатике}{}{}{}{
\begin{itemize}
\item Диплом I степени.
\end{itemize}}

\cventry{2017}{Всероссийская командная олимпиада школьников по программированию (ВКОШП)}{}{}{}{
\begin{itemize}
\item Призер
\end{itemize}}


%--------------------------------------
%	COMMUNICATION SKILLS SECTION
%--------------------------------------


\section{Extra Curriculars}

\cventry{Май 2018 -- настоящее время}{Репетитор по олимпиадному программированию}{}{}{}{Мои ученики успешно сдают вступительные экзамены в математические классы 57 и 179 школы г. Москвы. Некоторые из них также поступают на кружки по программированию от Tinkoff и Яндекс. На данный момент шестеро моих учеников учатся в ведущих математических школах страны.
}


%--------------------------------------
%	REFERENCES SECTION
%--------------------------------------


\section{References}

\begin{multicols}{2}

\cventry{}{Ященко Иван Валериевич}{\newlineПреподаватель НИУ ВШЭ,\newlineПреподаватель школы №57,\newlineДиректор МЦНМО}{}{\newline +7(916)810-34-35}{}
\columnbreak

\cventry{}{Чумаков Олег Олегович}{\newline CEO Luden.io}{}{\newline oleg.chumakov@luden.io}{ }
\end{multicols}

\end{document}